\chapter{Теория}

\section{Перечень организаций по стандартизации}

\subsection{ISOC (Internet Society - общество Интернета)}
Общество Интернета (англ. Internet Society, ISOC) — международная профессиональная
организация, занимающаяся развитием и обеспечением доступности сети Интернет.
Организация насчитывает более 20 тысяч индивидуальных членов и более 100 организацийчленов в 180 странах мира. Общество Интернета предоставляет организационную основу для
множества других консультативных и исследовательских групп, занимающихся развитием
Интернета, включая IETF и IAB.

\subsection{IRTF (Internet Research Task Force - целевая группа интернетисследований)}

IRTF (англ. Internet Research Task Force — Исследовательская группа Интернеттехнологий, www.irtf.org) — подразделение IAB, которое выполняет долгосрочные
исследовательские программы, связанные с вопросами развития архитектуры, базовых
протоколов и сетевых приложений сети Интернет.

\subsection{IETF (Internet Engineering Task Force - рабочая группа по инженерным вопросам интернета)}
Инженерный совет Интернета (англ. Internet Engineering Task Force, IETF) — открытое
международное сообщество проектировщиков, учёных, сетевых операторов и провайдеров,
созданное IAB в 1986 году и занимающееся развитием протоколов и архитектуры Интернета.

\subsection{IEEE (Institute of Electrical and Electronics Engineers - Институт инженеров электротехники и электроники)}
IEEE (I triple E — «Ай трипл и») — некоммерческая инженерная ассоциация из США,
разрабатывающая широко применяемые в мире стандарты по радиоэлектронике,
электротехнике и аппаратному обеспечению вычислительных систем и сетей.

\subsection{ICANN (Internet Corporation for Assigned Names and Numbers - Корпорация по управлению доменными именами и IP-адресами)}
ICANN (читается «айкэн») — международная некоммерческая организация, созданная 18
сентября 1998 года при участии правительства США для регулирования вопросов, связанных с
доменными именами, IP-адресами и прочими аспектами функционирования Интернета.
С 1 октября 2016 года — независимая международная организация.

\subsection{EIA (Electronic Industries Alliance - альянс электронной промышленности)}
Альянс отраслей электронной промышленности. Расположенная в США профессиональная организация, разрабатывает электрические и функциональные стандарты с
идентификатором RS (Recommended Standards).

\subsection{ITU (International Telecommunication Union - международный союз электросвязи)}
ITU - международная организация, определяющая рекомендации в области телекоммуникаций и радио, а также регулирующая вопросы международного использования
радиочастот (распределение радиочастот по назначениям и по странам). Основан как
Международный телеграфный союз в 1865 году, с 1947 года является специализированным
учреждением ООН.


Результатом
работы многих вышеперечисленных организаций являются такие документы, как: стандарты,
технические спецификации, положения о применимости и RFC (Request for Comments).
\endinput